\chapter{Diseño e implementación del modelo a desarrollar}
\label{cap:diseno}

Este capítulo articula el modelo de investigación y su materialización en software. Se presenta primero el \emph{diseño metodológico} —el tipo de investigación y el modelo de simulación numérica que la sostiene, las variables y su matriz de consistencia, los casos de estudio y su criterio de selección, el procedimiento de verificación y validación con sus instrumentos, y los criterios de aceptación con que se contrasta la hipótesis— y, a continuación, el software EduFEM, que concreta ese modelo.

\section{Diseño metodológico}
\label{sec:metodologia}

La Introducción presentó una visión general del enfoque de investigación; esta sección lo desarrolla siguiendo la estructura de un modelo de investigación: la conceptualización y el alcance del modelo, la definición de las variables, el procedimiento y la toma de observaciones, la validación de esas observaciones y el contraste de la hipótesis con los resultados.

\subsection{Tipo de investigación y modelo de simulación numérica}
\label{sec:proceso-met}

Por su finalidad, el trabajo es de carácter \emph{propositivo} —la propuesta de un artefacto que resuelve un problema formativo—; por su naturaleza es investigación aplicada de tipo \emph{tecnológico}, la que produce conocimiento operativo para transformar una realidad concreta \autocite[pp.~90-92]{garciacordoba2005tecnologica}; por el abordaje del problema, el enfoque es \emph{cuantitativo} en su fase de verificación y validación \autocite[p.~6]{sampieri2018metodologia}, pues la calidad del artefacto se establece midiendo errores numéricos, tensiones, desplazamientos y tasas de convergencia frente a referencias de exactitud conocida; y por su nivel es \emph{descriptiva} y \emph{comparativa}: caracteriza el comportamiento numérico del modelo y contrasta de forma sistemática el desempeño de los elementos Q4 y Q9 y el del motor frente a las referencias. El modelo que sostiene la investigación es un \emph{modelo de simulación numérica} —una representación ideal del objeto en la que el investigador abstrae y sistematiza las relaciones que considera esenciales \autocite[p.~21]{alvarez_metodologia}, ejercitada aquí por simulación computacional y no por experimentación física—: el objeto de estudio declarado en la Introducción —el análisis de medios continuos en elasticidad lineal plana por el MEF— se estudia resolviendo problemas de ese dominio con el motor de cálculo desarrollado y contrastando sus resultados con referencias de exactitud conocida; el campo de acción —la enseñanza de ese análisis— recibe el aporte en forma de los atributos del artefacto, que se cuentan con los indicadores de la \autoref{tab:variables}. No hay unidades experimentales, tratamientos ni aleatorización: el motor es determinista, de modo que la confiabilidad del modelo no descansa en réplicas sino en la reproducibilidad exacta de cada corrida y en la convergencia de sus resultados bajo refinamiento de malla; la incertidumbre numérica se cuantifica con las tasas de convergencia y, cuando la referencia no es exacta, con la extrapolación de Richardson de la propia secuencia de mallas \autocite[pp.~309-312]{oberkampf2010vv}.

El trabajo combinó dos ciclos articulados. El \emph{desarrollo del software} siguió un proceso iterativo e incremental: cada componente del canal de cálculo del MEF (\autoref{cap:marco_teorico}) se implementó y se sometió a prueba de regresión antes de integrarse al conjunto. La \emph{evaluación de la corrección} se realizó mediante un esquema de verificación y validación (\autoref{sec:verificacion-validacion}): verificación de código por el método de soluciones manufacturadas y contrastación con dos referencias externas, la solución analítica de la viga de Timoshenko y la membrana de Cook \autocite{oberkampf2010vv,salari2000mms}.

\subsection{Variables e indicadores}
\label{sec:variables-met}

La \autoref{tab:variables} operacionaliza las variables del estudio, es decir, especifica para cada una las operaciones concretas con que se la mide \autocite[p.~137]{sampieri2018metodologia}. Las variables del problema científico —la forma en que la herramienta expone el canal de cálculo y la observabilidad y contrastabilidad del procedimiento— se operacionalizan como \emph{atributos del artefacto} con indicadores contables: etapas del canal con módulo interactivo y con desarrollo numérico en la memoria, fases que comparten el lienzo y formatos con prueba de ida y vuelta. Las variables del experimento numérico con que se establece la corrección del motor son la \emph{variable independiente} —la definición del modelo estructural (geometría, material, acciones y apoyos) y las decisiones de discretización: el tipo de elemento (Q4 o Q9) y la densidad de malla (tamaño característico $h$, o $N\times N$)— y la \emph{variable dependiente} —la respuesta estructural calculada (tensiones y desplazamientos) junto con su \emph{exactitud} frente a las soluciones de referencia, medida por las normas de error $L^2$ y $H^1$ del desplazamiento, por la norma $L^2$ del campo de tensiones recuperado y por los errores relativos en tensión normal y en flecha—. El tipo de análisis (tensión o deformación plana) y el orden de cuadratura se mantienen como \emph{variables controladas}, y el desempeño del solucionador se registra como \emph{variable observada}, sin criterio de aceptación, para documentar la escalabilidad del motor.

\begin{table}[htbp]
\centering
\caption{Operacionalización de las variables de la investigación.}
\label{tab:variables}
\begin{tabular}{@{}p{2.7cm}p{3.2cm}p{5.1cm}p{2.5cm}@{}}
\toprule
\textbf{Variable} & \textbf{Dimensión} & \textbf{Indicador} & \textbf{Instrumento} \\
\midrule
Atributos del artefacto (variables del problema científico) & Cobertura del canal; organización de la interfaz; trazabilidad e interoperabilidad & Etapas del canal con módulo interactivo y con desarrollo numérico en la memoria; fases que comparten el lienzo; formatos con prueba de ida y vuelta & Inventario de módulos; memoria del ejemplo canónico; pruebas de interoperabilidad y de generación de la memoria \\
Independiente: definición del modelo & Geometría; material; acciones y apoyos & Dominio y malla; $E$, $\nu$, espesor $t$; cargas y restricciones & Guion de cada caso (malla, material, cargas y apoyos) \\
Independiente: discretización & Tipo de elemento; densidad de malla & Q4/Q9; tamaño característico $h$ (o $N\times N$) y número de GDL & Guion de cada caso (tipo de elemento y $N$) \\
Dependiente: respuesta estructural & Tensiones; desplazamientos & $\sigma_x,\sigma_y,\tau_{xy},\sigma_{\mathrm{VM}}$; $u,v$ nodales & Solucionador del MEF \\
Dependiente: exactitud de la solución & Convergencia; error en desplazamiento; error en tensión; equilibrio & Tasa de convergencia (orden $p$); normas $L^2$ y $H^1$ del desplazamiento y $L^2$ de $\sigma^*$; error relativo (\%) en $\sigma_x$ y en flecha frente al analítico y a SAP2000; residuo de la suma de reacciones & Guiones de V\&V \\
Controladas & Tipo de análisis; integración & Tensión o deformación plana; orden de cuadratura & Constantes del guion de cada caso \\
Observada: desempeño & Tiempo; memoria & Tiempos de ensamblaje y de solución; memoria de $\bm{K}$ densa y dispersa & Guion de medición de tiempos \\
\bottomrule
\end{tabular}
\end{table}

\subsection{Matriz de consistencia}
\label{sec:matriz-consistencia}

El \textbf{problema científico}, en la formulación única de la Introducción —¿cómo debe construirse un software educativo para que exponga de forma transparente, interactiva y numéricamente verificable el canal de cálculo completo del MEF en elasticidad plana, de modo que cada etapa del procedimiento resulte observable y contrastable por el estudiante de ingeniería como apoyo a la comprensión de sus fundamentos?—, se atiende mediante el \textbf{objetivo general} de desarrollar EduFEM y se contrasta con la \textbf{hipótesis de diseño} de que un software con esos atributos reproducirá las tasas teóricas de convergencia con errores acotados, evidenciará el bloqueo por cortante y hará observable cada etapa del canal. La \autoref{tab:consistencia} traza la correspondencia entre cada objetivo específico, la pregunta de investigación que responde, las variables o indicadores asociados, el método o instrumento empleado y la evidencia que lo respalda en el documento.

\begin{table}[htbp]
\centering
\caption[Matriz de consistencia de la investigación.]{Matriz de consistencia: correspondencia entre objetivos específicos, preguntas de investigación, variables, métodos y evidencia.}
\label{tab:consistencia}
{\footnotesize
\setlength{\tabcolsep}{4pt}
\renewcommand{\arraystretch}{1.2}
\begin{tabular}{@{}p{3.0cm}p{2.3cm}p{2.7cm}p{2.7cm}p{2.4cm}@{}}
\toprule
\textbf{Objetivo específico} & \textbf{Pregunta asociada} & \textbf{Variable / indicador} & \textbf{Método / instrumento} & \textbf{Evidencia} \\
\midrule
Fundamentar teóricamente el MEF en elasticidad plana & PI-1 a PI-3 & Canal de cálculo cubierto & Revisión de la literatura clásica del MEF y de la V\&V & \autoref{cap:marco_teorico} \\
Implementar el motor MEF 2D Q4/Q9 & PI-1 & Tipo de elemento; exactitud (normas $L^2$, $H^1$ y $L^2$ de $\sigma^*$) & Implementación pura en NumPy/SciPy; pruebas de regresión & \S\ref{sec:motor}; \autoref{cap:resultados} (MMS) \\
Diseñar la interfaz pre/proceso/post & PI-2 & Fases que comparten el lienzo; sincronización lienzo--tabla & Diseño centrado en un lienzo único compartido & \S\ref{sec:pre-proceso}--\S\ref{sec:post-proceso}; \autoref{tab:fases} \\
Desarrollar los módulos educativos & PI-2 & Etapas del canal con módulo interactivo & Capas interactivas superpuestas a la malla real; inventario de módulos & \S\ref{sec:educacion}; \autoref{sec:cobertura} \\
Verificar y validar el software & PI-1, PI-3 & Exactitud; tasa de convergencia; bloqueo por cortante & MMS, viga de Timoshenko, membrana de Cook & \autoref{cap:resultados} \\
Generar la memoria de cálculo e interoperabilidad & PI-2 & Etapas con desarrollo numérico en la memoria; formatos con ida y vuelta & Generación de PDF con las funciones del motor; importación/exportación DXF y CSV & \S\ref{sec:memoria-io}; \autoref{anexo:memoria}; \autoref{sec:consistencia-interna} \\
\bottomrule
\end{tabular}}
\renewcommand{\arraystretch}{1.0}
\end{table}

Las preguntas de investigación referidas en la matriz son las enunciadas en la Introducción: \textbf{PI-1}, sobre la precisión verificable del motor frente a soluciones analíticas y comerciales; \textbf{PI-2}, sobre la organización de la interfaz y de los módulos que expone el canal de cálculo de forma transparente; y \textbf{PI-3}, sobre la evidencia de fenómenos numéricos como el bloqueo por cortante del Q4 frente a la convergencia del Q9.

El primer objetivo específico —la fundamentación teórica— es transversal: sustenta las tres preguntas de investigación y se materializa en el \autoref{cap:marco_teorico}, de ahí que la matriz le asocie las tres. El sexto objetivo específico responde a \textbf{PI-2} en su parte principal: la memoria de cálculo es la forma escrita en que el canal se expone de manera transparente y contrastable, y su cobertura se cuenta con el mismo indicador que la de los módulos. La interoperabilidad de datos (DXF, CSV) es su complemento instrumental: no deriva de una pregunta de investigación, pero se verifica por ida y vuelta e idempotencia (\autoref{sec:consistencia-interna}) para que el modelo del alumno pueda entrar y salir de la herramienta sin pérdida.

\subsection{Casos de estudio y criterio de selección}
\label{sec:poblacion}

En el esquema clásico, la población sería el universo de problemas de elasticidad lineal plana —en tensión o deformación plana— resolubles con elementos cuadriláteros isoparamétricos. El conjunto de casos que sigue no es una muestra estadística de ese universo, sino una selección \emph{intencional} por valor probatorio —una muestra dirigida, no probabilística \autocite[p.~215]{sampieri2018metodologia}—, coherente con la práctica de la verificación de códigos de cálculo, donde los casos de prueba se eligen por su capacidad de ejercitar todos los términos del modelo matemático y por cubrir la topología de malla más general que el código deberá resolver, antes que por su realismo físico \autocite{oberkampf2010vv}. Cada caso aporta una referencia de naturaleza distinta y ejercita términos distintos del modelo. El método de soluciones manufacturadas provee una solución exacta arbitraria y ejercita la fuerza de volumen, las restricciones de Dirichlet no homogéneas, las dos matrices constitutivas y el mapeo sobre elementos distorsionados, en cuatro configuraciones (cuadrado unitario y cuadrilátero distorsionado, en tensión y en deformación plana). La viga de Timoshenko simplemente apoyada ejercita la carga superficial uniforme y los apoyos, y aporta las dos referencias externas: una solución analítica de la teoría de la elasticidad y un modelo de cáscara en SAP2000. La membrana de Cook ejercita la distorsión trapezoidal bajo cortante y flexión, el modo de deformación ante el que el Q4 bloquea. Entre los tres quedan cubiertos los términos del modelo de la \autoref{sec:formulacion-debil} —cargas nodales, volumétricas y superficiales uniformes, restricciones homogéneas y no homogéneas, ambos estados planos y ambos elementos—, con una excepción que se declara: la carga superficial linealmente variable no interviene en ningún caso de referencia y no tiene contraste externo; las recomendaciones proponen un caso que la incorpore. El refinamiento de malla y la comparación Q4 frente a Q9 se realizan en los dos casos con secuencia de mallas —el MMS y la membrana de Cook, ambos con $N\in\{2,4,8,16,32\}$—; la viga de Timoshenko se resuelve con una única malla fina de elementos Q9 y constituye un contraste puntual de exactitud frente a las dos referencias externas, no un estudio de convergencia.

\subsection{Procedimiento, instrumentos y reproducibilidad}
\label{sec:validacion-datos}

Cada caso sigue el mismo procedimiento. Un guion construye la malla y define material, cargas y apoyos; un validador de salud del modelo —una función pura sin dependencias de la interfaz— verifica la consistencia de los datos de entrada antes de resolver: existencia de elementos, restricciones suficientes para suprimir los movimientos de cuerpo rígido, ausencia de nodos referenciados inexistentes y no degeneración de los elementos (determinante Jacobiano positivo); los errores críticos bloquean el cálculo. En la verificación por soluciones manufacturadas la consistencia de los datos está garantizada por construcción: el término fuente se deduce analíticamente de la solución exacta elegida a priori, de modo que las cargas y las restricciones de Dirichlet corresponden exactamente a ella. El sistema se ensambla y resuelve con las mismas funciones del motor que consumen los módulos educativos y la memoria de cálculo, lo que garantiza coherencia entre lo que el alumno ve explicado y lo que el programa resuelve, y las tensiones se recuperan por la secuencia descrita en el \autoref{cap:marco_teorico}.

Los instrumentos de medición son los guiones de verificación y validación del proyecto (\texttt{tests/vv\_mms.py}, \texttt{tests/vv\_timoshenko.py} y \texttt{tests/vv\_cook.py}), que computan las normas de error —integradas con un punto de Gauss más por dirección que la rigidez (\autoref{sec:verificacion-validacion}) para no subestimar artificialmente el error—, los errores relativos frente a las referencias y las tasas de convergencia, y depositan los datos crudos en archivos CSV. Dentro de cada caso solo se manipulan el tipo de elemento y la densidad de malla —fijos en la viga de Timoshenko, por lo dicho en la \autoref{sec:poblacion}—; el tipo de análisis, el material, la geometría, las cargas y los apoyos, el orden de cuadratura ($2\times2$ para Q4 y $3\times3$ para Q9) y las tolerancias numéricas del motor —el cero numérico y el determinante Jacobiano mínimo admisible, constantes del programa— permanecen fijos.

La reproducibilidad se apoya en el determinismo del motor: ante las mismas entradas produce las mismas salidas, y cada caso se regenera por completo al reejecutar su guion con las dependencias declaradas en el repositorio. A ello se suma una batería de pruebas de regresión que verifica la estabilidad del modelo de datos ante operaciones de edición —el ciclo de transformación $\text{Q4}\rightarrow\text{Q9}\rightarrow\text{Q4}$ sin deriva numérica, la equivalencia de resultados con identificadores de nodo no contiguos y la interoperabilidad de formatos— y que se reporta en el \autoref{cap:resultados}.

\subsection{Criterios de aceptación y contraste de la hipótesis}
\label{sec:validacion-resultados}

Los resultados se contrastan con referencias de naturaleza distinta y con criterios fijados a priori en los propios guiones, que los evalúan automáticamente y terminan con error si alguno falla. Para el motor: tasas de convergencia del desplazamiento observadas dentro de $\pm0{,}5$ de las teóricas ($\mathcal{O}(h^{p+1})$ en $L^2$ y $\mathcal{O}(h^{p})$ en $H^1$) en las cuatro configuraciones del MMS; para el campo de tensiones recuperado, cuya tasa teórica no está establecida para la secuencia de extrapolación y promediado nodal que emplea el motor, una cota empírica mínima declarada como tal —tasa no inferior a $1{,}4$ en Q4 ni a $1{,}9$ en Q9—; en la viga de Timoshenko, error de la flecha central inferior al $3\,\%$ frente a la solución analítica con corrección de cortante, error de la tensión normal $\sigma_x$ inferior al $1\,\%$ frente a la solución analítica y frente a SAP2000 en los tres puntos de control, y equilibrio global —suma de reacciones verticales igual a la carga total— con residuo relativo inferior a $10^{-8}$; y, en la membrana de Cook, desplazamiento del Q9 con $N=8$ dentro del $1{,}5\,\%$ del valor de referencia y flecha del Q4 inferior a la del Q9 con la misma malla, que es la firma del bloqueo por cortante. Para los atributos del artefacto: módulo interactivo para cada una de las siete etapas elementales y de ensamblaje del canal (M1 a M7) y desarrollo con sustitución numérica de las nueve etapas en la memoria de cálculo; las tres fases del análisis operando sobre un mismo lienzo; ida y vuelta CSV/ZIP sin pérdida de entidades ni de resultados, reimportación DXF idempotente y generación de la memoria para Q4 y Q9 sin error. El desempeño del solucionador se reporta sin criterio de aceptación. El \autoref{cap:resultados} reporta los valores efectivamente medidos en cada magnitud, incluidas las componentes secundarias de tensión que no forman parte del criterio.

La hipótesis de diseño enunciada en la Introducción se contrasta cláusula por cláusula con esos criterios: la cláusula (a) —tasas teóricas y errores acotados frente a Timoshenko y SAP2000— y la cláusula (b) —bloqueo del Q4 frente a la convergencia del Q9— con los criterios numéricos del motor (PI-1 y PI-3), y la cláusula (c) —cada etapa del canal observable en un módulo o en la memoria— con los indicadores de cobertura de los atributos del artefacto (PI-2). En la nomenclatura del diseño clásico de investigación, el \emph{aporte} es EduFEM; el \emph{atributo manipulado} —la transparencia, interactividad y verificabilidad con que se expone el canal de cálculo, frente a la opacidad del software profesional y la abstracción de los textos clásicos— se materializa en las decisiones de diseño de la \autoref{sec:diseno-edufem} y se cuenta con los indicadores de cobertura, fases y formatos de la \autoref{tab:variables}; y el \emph{efecto esperado} —la observabilidad y contrastabilidad del procedimiento— se establece cuando cada etapa del canal tiene una exposición interactiva o documental y cuando el valor que esa exposición muestra coincide con el que calcula el motor verificado, coincidencia que garantiza el uso de las mismas funciones en el motor, los módulos y la memoria (\autoref{sec:validacion-datos}). Estos tres elementos no deben confundirse con las variables del experimento numérico de la \autoref{tab:variables}, que miden la corrección del motor. El supuesto de que esa observabilidad apoya el aprendizaje no se contrasta: la evaluación empírica del impacto educativo con estudiantes excede el alcance temporal del trabajo y se plantea como línea futura.
